\scanpage{122}
% archive sheet number: 121

$\mathbb{X}$ \deleted{et}\inserted{\handwritten{(}donc} $X$ est régulier en les points de son support\handwritten{)},\deleted{$U=f^{\handwritten{-1}}(V)=X-\operatorname{supp}D$}\handwritten{.}
\margin{\handwritten{(Condition de transversalité)}}

\handwritten{b}) $\operatorname{supp}E$ est irréductible, et il est régulier en les points de $f(\operatorname{supp}D)=\operatorname{supp}E\handwritten{\cap}f(X)$ (donc $Y$ est régulier en ses points).

Alors pour tout $x\handwritten{\in}\operatorname{supp}D$, \deleted{\unsure{l'image des}}\inserted{les} sous-groupes d'inertie de $G$ en $x$ \unsure{ont|sont} comme images par $\handwritten{\phi}:G\handwritten{\longrightarrow}H$ des sous-groupes de $H$ qui sont des sous-groupes d'inertie de $H$ relativement au point maximal de $\operatorname{supp}E$. Par suite, pour toute représentation linéaire de $G$ (dans un module $M$ sur quelque anneau) \deleted{\unsure{sur}} qui veut bien se factoriser par une représentation linéaire \uline{de} $H$, les images dans $\operatorname{Aut}(M)$ des \inserted{sous-}groupes d'inertie \deleted{\unsure{dans}} de $G$ relatifs aux divers points de $\operatorname{supp}D$ (ou encore, relatifs aux diverses composantes irréductibles de $\operatorname{supp}D$, cela revient au même en vertu de 3\handwritten{)} sont conjugués dans $\operatorname{Aut}(M)$.

\inserted{\uline{Remords\handwritten{\unsure{:}} au n\textsuperscript{o}}\handwritten{s3,4}.} On peut demander \deleted{\unsure{un}}\handwritten{\unsure{des|du}} résultat\handwritten{\unsure{s}} plus précis, en \deleted{\uline{Application}} considérant pour tout $\handwritten{\unsure{\xi|\zeta}}$ comme ci-dessus (au dessus d'un point géométrique\inserted{$\handwritten{\bar}{x}$} de $\operatorname{supp}D$) l'homomorphisme correspondant $\underline{Z}_{\handwritten{\unsure{\ell|l}}}(1)(\handwritten{\bar}{x})\handwritten{\longrightarrow}G$ dont l'image est le sous-groupe d'inertie correspondant, et en se demandant comment cet homomorphisme varie avec le choix de $\handwritten{\xi}$ \uline{et de} $\handwritten{\bar}{x}$. Pour $\handwritten{\xi}$ variable, la variation se fait évidemment modulo composition avec les automorphismes intérieurs de $G$, \deleted{Pour donner un sens à la question} donc pour $\handwritten{\bar}{x}$ fixé, on trouve une classe de conjugaison d'homomorphismes
\[
\underline{Z}_{\handwritten{\ell}}(1)(\handwritten{\bar}{x})\handwritten{\longrightarrow}G\ .
\]
\deleted{pour}\inserted{\handwritten{Étudions}} \deleted{L}\handwritten{l}a variation de cet homomorphisme\handwritten{,} pour $\handwritten{\bar}{x}$ variable dans une composante \deleted{de $D$}\inserted{$D_i$ \handwritten{de}} l'ensemble des points réguliers de $D$\inserted{à car. $\handwritten{\ne\ell}$}; d'après ce qu'on a vu au n\textsuperscript{o}\unsure{3|2}, on passe d'un tel homomorphisme à un autre en composant avec un quelconque des isomorphismes $\underline{Z}_{\handwritten{\ell}}(1)(\handwritten{\bar}{x})\handwritten{\longrightarrow}\underline{\unsure{\mathbb{Z}|E|Z}}_{\handwritten{\ell}}(1)(\handwritten{\overline}{x'})$ associé à une classe de chemins de $\handwritten{\bar}{x}$ à $\handwritten{\overline}{x'}$. Si cet isomorphe ne dépend pas du choix de cette classe de chemins, i.e. si le\deleted{action du groupe fondamental} faisceau \unsure{l|1}-adique de \handwritten{T}ate $\underline{Z}_{\handwritten{\ell}}(1)$ sur $D(\handwritten{\ell})$ est trivial (ce qui se voit déjà sur la

% Contour separating the Remords passage from the preceding paragraph; doubled left boundary.
\figure{\handwritten{\begin{tikzpicture}[x=0.6cm,y=0.6cm]
\draw (15,12.5) .. controls (12,12.2) and (10,12.4) .. (5,12.4) -- (4.7,11.5) -- (3.3,11.5) .. controls (1.8,11.1) and (1.4,11.4) .. (0,11.3) .. controls (-0.2,10) and (0.2,7) .. (-0.1,6) .. controls (-0.3,4) and (0.3,2.3) .. (0.1,2) .. controls (-0.2,1) and (0,0.4) .. (0.2,0);
\draw (0,11.3) .. controls (-0.1,8.5) and (-0.2,8) .. (-0.1,6) .. controls (0,3.8) and (0,1.5) .. (0.25,0.4);
\end{tikzpicture}}}
